
%WDH Mealy und Moore Definition->Unterschied, Verwandschaft
\section{Endliche Automaten} %%%%%%%%%%%%%%%%%%%%%%%%%%%%%% Endliche Automaten %%%%%%%%%%%%%%%%%%%%
\subsection*{}
\begin{frame}
  \frametitle{Arten von Automaten}

  Es gibt zwei Arten, wie ein Automat eine Ausgabe tätigen kann. Wir unterscheiden dabei:
  \pause
  \begin{block}{Mealy-Automat}
    \begin{itemize}
      \item Erzeugung einer Ausgabe bei jedem Zustandsübergang
      \item Ausgabefunktion $g: Z \times X \rightarrow Y^*$
      \item Markieren der Kanten mit $x_i|y_i$
    \end{itemize}
   \end{block}
    \pause
   \begin{block}{Moore-Automat}
    \begin{itemize}
      \item Erzeugung einer Ausgabe bei Erreichen eines Zustands
      \item Ausgabefunktion $h: Z \rightarrow Y^*$
      \item Markieren der Zustände mit $q_i|y_i$ ($q_i$ ist Zustandsname)
    \end{itemize}
  \end{block}

  In beiden Fällen ist die Ausgabe ein Wort $y = y_0\ldots y_{n-1}$ über einem
  Ausgabealphabet $Y$.
\end{frame}


\begin{frame}
   \frametitle{Der Akzeptor}

   \begin{itemize}
     \item \vspace{-1.4em} Ist der häufigste \textbf{Spezialfall} eines Moore-Automaten
     \item Eine Ausgabe findet nicht bei allen Zuständen statt \pause
     \item Die Zustände $F \subseteq Z$, bei denen eine Ausgabe (immer ein Bit lang) erfolgt, heißen \texttt{akzeptierende Zustände} \\
     Es gilt $F=\{z|h(z)=1\}$ \pause
     \item graphisch werden diese durch Doppelkreise angegeben
     		\begin{center}
        \begin{tikzpicture}
    			\draw(0,0) circle(1.4ex);
      		\node[draw,circle] (qi) at (0,0) {$z_f$};
      	\end{tikzpicture}
      	\end{center} \pause
     \item[]
     \item Ein Wort $w\in X^*$ wird \texttt{akzeptiert}, wenn gilt $f^*(z_0,w)\in F$
     \item Die von einem Akzeptor $A$ \texttt{akzeptierte formale Sprache} ist $L(A)=\{w \in X^* |f^*(z_0,w)\in F\}$
   \end{itemize}
 \end{frame}

%%%%%%%%%%%%%%%%%%%Reguläre Ausdrücke%%%%%%%%%%%%%%%%%%%%%%%%%
\section{Reguläre Ausdrücke}
\subsection*{}

\begin{frame}
  \frametitle{Definition}

  Reguläre Ausdrücke sind eine verbreitete und geeignete Notation, um reguläre
  Sprachen (Typ-3) zu formalisieren.

	\begin{block}{Die Regeln (= Metazeichen)}
		\begin{tabular}{c p{0.7\textwidth}}
				\xb{Metazeichen} & 	\xb{Bedeutung} \\
				$( )$			& Klammerung von Alternativen\\
				$*$			& n-maliges Vorkommen\\
				$|$			& trennt Alternativen\\
		\end{tabular}
	\end{block}

  \begin{block}{Es gelten folgende Vorrangregeln:}
   \begin{itemize}
      \item $*$ bindet stärker als Verkettung
      \item Verkettung $(RS)$ bindet stärker als "`oder"' $(R|S)$
      \item Überflüssige Klammern dürfen wir weglassen. \\
      So sind $(RS),((RS)),\ldots$ und $RS$ äquivalent
    \end{itemize}
  \end{block}
\end{frame}


\subsection*{}
\begin{frame}
  \frametitle{Definition}
	\begin{block}{Die Sprache von R}
		Wenn \textbf{R} ein regulärer Ausdruck ist, dann bezeichnen wir mit \textbf{$\langle R \rangle$} die Sprache, die dieser erzeugt.\\
  		\begin{itemize}
    			\item $\langle \emptyset \rangle =\{\}$
			\item Für $a \in A$ ist $\langle a \rangle=\{a\}$
			\item $\langle R_1 | R_2 \rangle = \langle R_1 \rangle \cup \langle R_2 \rangle$
			\item $\langle R_1 R_2 \rangle = \langle R_1 \rangle \cdot \langle R_2 \rangle$
  		\end{itemize}
	$R_1$ und $R_2$ sind hier zwei beliebige reguläre Ausdrücke.


	\end{block}
\end{frame}

\subsection*{}
\begin{frame}
  \frametitle{Die Sprache von R}
	\begin{block}{Beispiel 1}
		Welche Wörter erzeugt der folgende reguläre Ausdruck R?
  		\begin{itemize}
			\item $R=(a|b)*abb(a|b)*$ ? \pause
			\item $\langle R \rangle$ enthält genau die Wörter, in denen das Teilwort $abb$ vorkommt.
  		\end{itemize}
	\end{block}

	
	\begin{block}{Beispiel 2}
		Gebe einen regulären Ausdruck für die Sprache aller Wörter die nicht $ab$
		enthalten
  		\begin{itemize} \pause
			\item $ {b*} {a*}$
  		\end{itemize}
	\end{block}
	
\end{frame}

\subsection*{}
\begin{frame}
  \frametitle{Die Sprache von R}
	\begin{block}{Beispiel 3}
		Welcher reguläre Ausdruck R erzeugt die Sprache $\{\epsilon\}$?
  		\begin{itemize}
			\visible<2->{\item $\emptyset*$, denn $\langle \emptyset \rangle ^* =
			\{\}^* = \{\epsilon\}$}
  		\end{itemize}
	\end{block}

	\visible<3->{
	\begin{block}{Beispiel 4}
		Gebe einen regulären Ausdruck für die Sprache aller Wörter mit mindestens 3 b's an!
  		\begin{itemize}
			\visible<4->{\item $(a|b)*b(a|b)*b(a|b)*b(a|b)*$ oder
			\visible<5->{\item $a*ba*ba*b(a|b)*$}}
  		\end{itemize}
	\end{block}
	}
\end{frame}

\subsection*{}
\begin{frame}
  \frametitle{Reguläre Aufgabe}

  Gegeben ist folgende Klasse von Sprachen über dem Alphabet $\Sigma = \{a,b,c\}$:
  \begin{multline*}
    L_n = \{w|
      \text{Das Wort w enthält genau einmal}\\
      \text{eine Folge von $a$ der Länge $n$, die nicht Teil}\\
      \text{einer Folge von $a$ mit einer Länge $>n$ ist}
          \}
  \end{multline*}

  Gebt einen regulären Ausdruck $R$ für $L_4$ an! (also: $\langle R \rangle = L_4$)

  \hide

  \begin{block}{Lösung a)}
    Wir stellen sicher, dass die vier $a$ genau einmal vorkommen, und sonst nur 1,2,3 oder mehr als 4 am Stück.
    \begin{multline*}
      R = ( (b|c)* (\emptyset |a|aa|aaa|aaaaaa*)(b|c)(b|c)*)* \\ aaaa (
      (b|c)(b|c)*(\emptyset |a|aa|aaa|aaaaaa* )(b|c)*)*
    \end{multline*}
  \end{block}
\end{frame}

\subsection*{}
\begin{frame}
  \frametitle{Ihr seid dran... Regulär Einkaufen}
  Gebt einen regulären Ausdruck an, der das Einkaufen in einem Supermarkt
  beschreibt. Beschreibt den Einkauf mit diesen Zeichen:

  \begin{description}
    \item[$b$] Laden betreten
    \item[$v$] Laden verlassen
    \item[$s$] Einkaufswagen verschieben
    \item[$p$] Produkt in den Einkaufswagen legen
    \item[$z$] Einkäufe bezahlen
  \end{description}

\only<1>{
  Den Einkaufswagen erhält man am Eingang beim Betreten des Ladens und gibt ihn
  beim Verlassen am Ausgang zurück. Die Produkte sind im ganzen Laden verteilt
  und nicht in Reichweite des Ein- oder Ausgangs -- aber an der Kasse gibt es Süßes!

  Beachtet, dass nichts zu bezahlen ist, wenn keine Waren im
  Einkaufswagen liegt (nur genau dann).

  Zielloses Rumstöbern ist erlaubt!
}
\only<2>{
  Eine mögliche Lösung ist:

  \begin{equation*}
    bss*(\emptyset | p(p|s)*zs)v
  \end{equation*}
}
\end{frame}

\subsection*{}
\begin{frame}
  \frametitle{von $\langle R \rangle$ zu $R$}
	\begin{block}{Und noch ein Beispiel...}
		Wenn $R$ ein regulärer Ausdruck für eine formale Sprache $L=\langle R \rangle$ ist, wie sieht dann ein regulärer Ausdruck
  		\begin{itemize}
			\item für $L^*$ aus? \visible<2->{  Lösung: $(R)*$}
			\item für $L^+$ aus? \visible<3->{  Lösung: $R(R)*$}
  		\end{itemize}
	\end{block}
\end{frame}

%%%%%%%%%%%%%%%%%%%Rechtslineare Grammatiken%%%%%%%%%%%%%%%%%%%%%%%%%
\section{Rechtslineare Grammatiken}

\subsection*{}
\begin{frame}
\frametitle{Beziehung zwischen kontextfreien und rechtslinearen Grammatiken}
	Wir können die Grammatiken in vier Klassen einteilen, Typ-0 bis Typ-3.\\
	Dabei gilt:
  	\begin{itemize}
    		\item Typ-$n$ $\supsetneq$ Typ-$(n+1)$
    		\item Je kleiner die Nummer, desto "`weniger eingeschränkt"' ist die Grammatik
    		\item Auch die Sprachen werden in diese Klassen gepackt: \\
    		Jede Sprache hat die Typ-Klasse der einfachsten Grammatik, die sie
    		erzeugt. \\ (Einfach entspricht höherer Typ-Nummer)
  	\end{itemize}
\end{frame}

\subsection*{}
\begin{frame}
	\frametitle{Übersicht über die Klassen}
	\begin{block}{Die verschiedenen Klassen}
  	\begin{description}
    		\item[$Typ$ $0$ $\&$ $1$] Typ-0 und Typ-1 kamen noch nicht vor
    		\item[$Typ$ $2$] Typ-2-Grammatiken sind die uns schon bekannten kontextfreien Grammatiken
    		\item[$Typ$ $3$] Typ-3-Grammatiken sind die neu hinzukommenden rechtslinearen Grammatiken
  	\end{description}
  	\end{block}
\end{frame}

\subsection*{}
\begin{frame}
	\frametitle{rechtslineare Grammatiken}
	\begin{block}{etwas genauer...}
  		Eine rechtslineare Grammatik ist eine kontextfreie Grammatik $G=(N,T,S,P)$ mit folgenden Einschränkungen. Jede Produktion ist entweder von der Form
  		\begin{itemize}
			\item $X \rightarrow w$ oder
			\item $X \rightarrow wY$ mit $w \in T^*$ und $X,Y \in N$
  		\end{itemize}
	\end{block} \pause

	\begin{block}{Regex}
    Zu jeder rechtslinearen Grammatik gibt es:
		\begin{itemize}
			\item \ldots einen entsprechenden regulären Ausdruck und \pause
			\item \ldots einen einen deterministischen endlichen Automaten
        \end{itemize}
    Zu jeder rechtslinearen gibt es äquivalente linkslineare Grammatiken. Diese
    "`können"' nichts anderes als rechtslineare Grammatiken, daher ignorieren
    wir sie in dieser Vorlesung.
	\end{block}
\end{frame}

\subsection*{}
\begin{frame}
	\frametitle{Semantik und Syntax regulärer Ausdrücke}
	\begin{block}{Vorsicht!}
		\begin{itemize}
			\item die Bedeutung, also die Semantik,  eines regulären Ausdrucks kann
			immer von einer rechtslinearen Grammatik angegeben werden \pause
			\item die Syntax regulärer Ausdrücke ist aber Typ-2 \pause
			\item d.h. eine Grammatik die alle regulären Ausdrücke zu einem Alphabet
			erzeugen kann, muss mindestens kontextfrei sein!
        \end{itemize}
  	\end{block}
\end{frame}

\subsection*{}
\begin{frame}
	\frametitle{rechtslineare Grammatiken}
	\begin{block}{Ein Beispiel \only<2->{- oder auch nicht...}}
  		Gegeben Sei die Grammatik $G=(\{X,Y\},\{a,b\},X,\{X \rightarrow aY | \epsilon, Y \rightarrow Xb\})$\\
  		\begin{itemize}
  		\item Ist diese Grammatik rechtslinear? \\
		\visible<2->{$G$ ist offensichtlich nicht rechtslinear, denn die Produktion $Y \rightarrow Xb$ hat das Nichtterminalsymbol links vom Terminalsymbol \\
		(Die Produktion ist linkslinear)!\\
		\visible<3->{\item Die Grammatik erzeugt die Sprache $L(G)=\{a^k b^k | k \in \mathbb N_0 \}$}\\}
		\visible<4->{\item Kann es eine rechtslineare Grammatik für diese Sprache
		geben? \\ Ist diese Sprache regulär? \\ } 
		\visible<5->{ Nein, ist sie nicht! }
		\end{itemize}
	\end{block}
\end{frame}

\subsection*{}

\begin{frame}
\frametitle{von G zu L(G)}
	\begin{block}{Aufgabe}
		Betrachte $G=(\{X,Y,Z\},\{a,b\},X,P)$ mit $P=\{X \rightarrow aX | bY | \epsilon, Y \rightarrow aX | bZ | \epsilon, Z \rightarrow aZ | bZ \}$
  		\begin{itemize}
			\item Was ist $L(G)$?
			\item Geben Sie einen endlichen Automaten an, der L(G) akzeptiert.
			\item Lässt sich diese Grammatik noch vereinfachen?
  		\end{itemize}
	\end{block}
	\visible<2->{
	\begin{block}{Lösung}
      		\begin{tikzpicture}[shorten >=1pt,node distance=2cm,auto,initial text=,>=stealth]
        		\node[state,initial,accepting]  (q_0)                       {$0$};
       			\node[state,accepting]          (q_1) [right of= q_0] {$1$};
        		\node[state]                    (q_2) [right of= q_1] {$2$};
        		\path[->] (q_0) edge [loop below]      node        {$a$} ()
        		edge [bend right] node [swap] {$b$} (q_1)
        		(q_1) edge              node        {$b$} (q_2)
        		edge [bend right] node [swap] {$a$} (q_0)
        		(q_2) edge [loop below] node        {$a,b$} ()
        		;
      		\end{tikzpicture}
	\end{block}
	}
\end{frame}

\begin{frame}
	\frametitle{Ist doch alles das Gleiche, oder?}
	\begin{block}{Gleiche Sprache - andere Grammatik}
	Folgende Grammatiken erzeugen die gleiche Sprache
  	\begin{itemize}
    		\item $G=(\{X,Y,Z\},\{a,b\},X,P)$ mit $P=\{X \rightarrow aX | bY | \epsilon, Y \rightarrow aX | bZ | \epsilon, Z \rightarrow aZ | bZ \}$
		\visible<2->{\item $G=(\{X,Y\},\{a,b\},X,P)$ mit $P=\{X \rightarrow a X | b Y | \epsilon, Y \rightarrow aX | \epsilon \}$
		\visible<3->{\item $G=(\{X\},\{a,b\},X,P)$ mit $P=\{X \rightarrow aX | baX |
		b | \epsilon \}$}}
  	\end{itemize}
  	\end{block}
\end{frame}


\section{Abschluss}
% Studis anzuregen darüber nachzudenken, ob sie wirklich alles wissen, ansonsten nachlesen oder fragen nachträglich stellen, dann kann in der nächsten Woche nochmal drauf eingegangen werden
\subsection*{}
\begin{frame}
	\frametitle{Zum Schluss...}
	\begin{block}{Was ihr nun wissen solltet!}
	\begin{itemize}
	  	\visible<2->{\item Worin unterscheiden sich Mealy- und Moore-Automaten?}
		\visible<3->{\item Welcher der beiden ist näher mit einem Akzeptoren verwandt?}
		\visible<4->{\item Wie setzt sich ein gültiger regulärer Ausdruck zusammen?}
		\visible<5->{\item In welchem Verhältnis stehen reguläre Ausdrücke, rechtslineare Grammatiken und endliche Automaten?}
		\visible<6->{\item Wie verhalten sich rechtslineare Grammatiken und kontextfreie Grammatiken zueinander?}
    \end{itemize}
   	\end{block}

	\visible<7->{
	\begin{block}{Ihr wisst was nicht?}
		Stellt \textbf{jetzt} Fragen!
	\end{block}}
\end{frame}